\section{Related work}

\paragraph{Sycophancy as a behaviour.}
% Establish: it is measured, it is widespread, it has a training-data explanation.
% Cite: sharma2023 (towards understanding sycophancy), perez2022 (model-written evals),
% wei2023 (simple synthetic data reduces sycophancy), NEW: any 2025-2026 follow-ups.
Prior work establishes that assistants trained with human preference feedback tend to agree
with users at the expense of accuracy, that the tendency scales with model size and with RLHF
training, and that preference data rewarding agreement is a plausible cause
\citep{sharma2023sycophancy, perez2022discovering}. Mitigations have been proposed at the data level
\citep{wei2023simple}. This literature measures the behaviour in aggregate; our contribution is
orthogonal, holding the pressure type as an experimental variable and asking where the
behaviour lives inside the model.

\paragraph{Linear directions and activation steering.}
% Cite: arditi2024 (refusal is mediated by a single direction), turner2023 (actadd),
% zou2023 (representation engineering), li2023 (inference-time intervention).
A body of work reports that high-level behaviours are mediated by low-dimensional directions
that can be read off activations and manipulated \citep{zou2023repe, turner2023actadd, li2023iti}. The
closest methodological precedent for this paper is the finding that refusal in
instruction-tuned models is mediated by a single direction whose ablation removes the
behaviour \citep{arditi2024refusal}, from which we take both the difference-in-means construction
and the all-layer ablation. Our result is a partial negative by comparison: the analogous
direction for capitulation does not exist usably below 7B, and where it does exist its
ablation removes about a quarter of the behaviour rather than most of it.

\paragraph{Probe validity.}
% Cite: hewitt2019 (control tasks), belinkov2022 (probing classifiers survey).
% This is where our gate design gets justified against prior critique.
The gap between what a probe can decode and what a model uses has been raised repeatedly,
with control tasks and selectivity proposed as correctives \citep{hewitt2019control, belinkov2022probing}.
Our validation protocol is in that tradition: the within-question matched-pair construction
is a control for question difficulty, the shuffled-label null is a control for the fitting
procedure, and the pushback positive control establishes that the pipeline can find a
direction when one exists. Section~\ref{sec:h6res} adds an empirical data point to this
discussion, in that a direction significant against a well-constructed null still had no
detectable causal role.

\paragraph{Emergence with scale.}
% Cite: wei2022 (emergent abilities), schaeffer2023 (mirage).
% IMPORTANT: engage with schaeffer2023 directly, since we claim a threshold.
Claims that capabilities appear abruptly at scale \citep{wei2022emergent} have been challenged on
the grounds that apparent discontinuities can be artifacts of discontinuous metrics
\citep{schaeffer2023mirage}. Our headline measure, AUROC, is continuous, and the discontinuity we
report is in whether a continuous quantity crosses a threshold fixed in advance rather than
in the quantity itself. The underlying values rise from 0.617 to 0.778 between 3B and 7B in
Qwen, which is a large jump but not a discontinuity in the metric. We therefore describe the
gate crossing as a threshold effect on a continuous measurement, and we note in
Section~\ref{sec:h1} that the second family climbs gradually over the same range.
